% !TEX root = ../thesis.tex

\chapter{Použitá metodológia}\label{ch:methodology}

Jednotlivé primitívy \emph{Robota Karla} budú vo vytváranej knižnici v jazyku \emph{C} implementované ako samostatné funkcie. Pre jej vytvorenie je však potrebné vybrať správnu a vhodnú knižnicu na vizualizáciu mikrosveta a navrhnúť formát mapy, ktorý bude použitý na reprezentáciu sveta \emph{Robota Karla}. Týmto témam bude venovaná táto kapitola.


\section{Výber knižnice pre vizualizáciu}

Pri rozhodovaní o tom, akú knižnicu pre vizualizáciu vybrať, bolo potrebné si ujasniť aj to, akým spôsobom budú študenti riešiť problémy \emph{Robota Karla}. V princípe sa jedná o dve možnosti:

\begin{enumerate}
   \item vytvorenie celého samostatného integrovaného prostredia na spôsob nástrojov od spoločnosti \emph{Borland} akým bol napríklad \emph{Turbo Pascal}, alebo

   \item vytvoriť len samostatnú knižnicu nezávislú od použitého vývojového prostredia.
\end{enumerate}

Aj napriek tomu, že má prvá možnosť svoje nesporné výhody a priamo sa ponúka otvorená implementácia rozhrania \emph{Turbo Vision} pomocou projektu \emph{tvision}\footnote{\url{https://github.com/magiblot/tvision}}, rozhodli sme sa vytvoriť len samostatnú knižnicu nezávislú od akéhokoľvek vývojového prostredia. Jedným z cieľov predmetu, v ktorom sa táto knižnica bude používať, je totiž aj naučiť sa základy používania Unix/Linux-ových textových editorov, akými je napríklad editor \emph{Vim}\footnote{\url{https://www.vim.org/}}.

\texttt{curses.h}\footnote{\url{https://invisible-island.net/ncurses/}}

\texttt{termbox}\footnote{\url{https://github.com/termbox/termbox2}}




\section{Vizualizácia mikrosveta}

vizualizacia sveta - inspiracia z clanku \cite{Untch1990}




\section{Formát sveta}


\section{Štruktúry pre vnútornú reprezentáciu}




% Syntetická časť opisuje metódy použité na syntézu riešenia a opisuje syntézu samotného riešenia (zvyčajne je to návrh/implementácia softvérového resp. hardvérového riešenia), pričom sa opiera o závery analytickej časti práce. Začína od toho, ako sa bude riešenie používať: najdôležitejšie scenáre používania a používateľské rozhranie, ktoré bude tieto scenáre efektívne podporovať. Až potom je na rade vnútorná architektúra alebo použité technológie. Syntetická časť tvorí zvyčajne ½ jadra práce.

% Syntetickú časť práce vhodne rozdeľte do kapitol a pomenujte ich podľa toho, čomu sú venované.
